\section{A complete resolution of the dimension-forty family}
All nine remaining rings of type $1^{\times4},2^{\times4},(\sqrt5)^{\times4}$
are now settled. Three are modeled, and six are excluded by a converse covering
every possible associator, rather than by a failed finite construction search.

\begin{lemma}\label{lem:cyclic-morita}
Let every noninvertible simple of an integral fusion category have dimension $n>1$
and the same left and right cyclic stabilizer $H$ of order $n$. Then its regular
subgroup algebra $A$ exists and its $A$-bimodule dual is pointed.
\end{lemma}
\begin{proof}
A fixed simple trivializes the pointed associator on $H$. Also, $H$ is normal,
since conjugation sends a common stabilizer to another such stabilizer.
Every noninvertible simple $X$ has $n$ simple right $A$-module structures because
$H^2(H,\mathbb C^*)=0$. For any one, $M$, its internal endomorphism object has
FP dimension $n^2/n=n$ and contains the unit once. Since every noninvertible has
dimension $n$, this internal end consists only of invertibles. Each constituent
fixes $X$, hence belongs to $H$, and occurs at most once. It is therefore exactly
$\bigoplus_{h\in H}h$. The cyclic action can be normalized to give a commuting
left $A$-action; its $n$ character choices yield $n^2$ simple bimodules over $X$,
each of normalized dimension one. An invertible source orbit has size $n$ and
cyclic diagonal stabilizer, again giving only normalized dimension-one bimodules.
This proves pointedness. The internal-end mechanism also appears in
\cite[Lemmas 6.1--6.2]{EGO}.
\end{proof}

In all nine source rings, the integral subcategory $\mathcal E$ has type
$1^{\times4},2^{\times4}$ and a common two-element stabilizer $\{1,g\}$ for its
dimension-two objects. The element $g$ is central in the entire based ring and
acts freely on the four irrational-dimensional objects. Thus $A=1+g$ exists,
and Lemma~\ref{lem:cyclic-morita} makes
$\mathcal D=\operatorname{Bimod}_{\mathcal C}(A)$ a category with twenty invertibles
and four $\sqrt5$-dimensional simples. It is a genuine generalized TY, or feudal,
category. Its noninvertible products are pointed. The reverse Morita algebra is
the actual internal end of the simple module $A$; its dimension is two, so it is
$1+c$ for an involution $c$ among the twenty invertibles.

Write $S=\Irr(\mathcal D_{\rm pt})$. In the reverse dual, the pointed integral
part contributes exactly $|N_S(\langle c\rangle)|$ invertibles. The source has four,
so this normalizer has order four. Among the five groups of order twenty, only
$F_{20}=C_5\rtimes_2 C_4$ and $D_{10}\times C_2$ have noncentral involutions.
For the latter group, the reverse dual integral ring is commutative for every
pointed associator admitting its subgroup algebra: all nonnormal $C_2$ double
cosets are inverse-stable with trivial stabilizer, so their dimension-two simples
are self-dual. The normalizer quotient is represented by the extra central
involution, giving equal involutive left/right invertible actions. Duality then
makes the remaining products commute. All eight noncommutative source rings have
noncommutative integral subrings. Their precursor must therefore have $S=F_{20}$.

The common stabilizer is the unique $C_5$ in $F_{20}$, and its associator is trivial
because it fixes a noninvertible simple. Hence the neutral category for the universal
grading is $\Vec_{C_5}$. Its Brauer--Picard group is
$O(C_5\oplus\widehat C_5,q)\cong D_8$ \cite{ENO}. The four invertible components
give the faithful order-four diagonal subgroup, through the multiplier two on
$C_5$; a rank-one noninvertible component gives a Fourier reflection. These
generate the full group. Thus the universal grading is $D_8$ and its homomorphism
to the Brauer--Picard group is an isomorphism. Changing the chosen noninvertible
component accounts for all reflection choices while preserving the distinguished
reverse involution $c=r^2$.

There is also a direct Liptrap-data proof of this restriction. If $t_m\in\mathbb F_5^*$
is the stabilizer bicharacter parameter, biderivation and quasisymmetry give
$t_{r^{-1}m}=2t_m$ and $t_{m^*}=t_m$. Writing the right quotient action as
an involution $\sigma$ of $C_4$ and the base noninvertible square as $p$ forces
$\sigma=-\mathrm{id}$ and $p=0$. The argument permits arbitrary nonzero complex
coherence variables and variable $\tau$. It is not a bounded phase-search obstruction.

The extension classification \cite{ENO} now applies to every hypothetical original
precursor. The product-system freedom is
$H^2(D_8,\operatorname{Inv}Z(\Vec_{C_5}))=H^2(D_8,C_5^2)=0$ by coprime averaging.
An explicit unitary precursor exists by its exact Liptrap data. Therefore all
associators form its $H^3(D_8,\mathbb C^*)$ torsor. This group has exponent dividing
four: restrict to an index-two $V_4$, whose third cohomology has exponent two,
then corestrict. Consequently every class has a fourth-root representative.
The map $H^3(C_2,\mu_4)\to H^3(C_2,\mathbb C^*)$ is an isomorphism, so the actual
reverse algebra condition is fully covered by the relative fourth-root classes
trivial on $\langle c\rangle$. Its multiplication is unique up to isomorphism,
since $H^2(C_2,\mathbb C^*)=0$.

The normalized bar complex has 343 three-cochain variables and 2402 constraints,
including the subgroup condition. Exact elimination gives a cocycle kernel of
order $2^{87}$ and a normalized coboundary image of order $2^{83}$: precisely
sixteen relative fourth-root classes. Independent algorithms check the cardinality
and all 120 representative-pair differences. These sixteen classes cover every
relative complex class, although different fourth-root classes can become equivalent
after extending coefficients.

For every representative the actual $1+c$ bimodules and balanced products are
constructed. Every resulting full tensor is exactly $\ring{156}$ or $\ring{263}$.
Two independent implementations recompute these sixteen outputs. The coherent
$C_4$ TY action separately realizes the commutative $\ring{449}$. Therefore
\[
 \boxed{\ring{309},\ \ring{383},\ \ring{467},\ \ring{470},\ \ring{501},\ \ring{514}
 \text{ are excluded over }\mathbb C.}
\]
This argument explicitly tracks the reverse algebra and all complex associator
classes; an arbitrary original category is never identified with one chosen unitary
model without the extension-classification proof.

\section{The dimension-twenty-four family}
Two exclusions follow directly from the general classification of invertible
pointed bimodule categories, with no assumption on the pointed associator.
Every component of a faithful fusion grading is an invertible bimodule category
over its neutral component \cite[Theorem 6.1]{ENO}. Its left and right stabilizers
must be isomorphic normal abelian groups: the mixed commutator is a nondegenerate
bicharacter between them \cite[Theorem 7.9 and Corollary 7.11]{GP}.

The supplied tensors for $\ring{445}$ and $\ring{493}$ have a faithful $C_3$
grading with pointed neutral component of rank eight. In a nontrivial component,
the exact stabilizers are
\begin{center}\begin{tabular}{lrrl}\toprule Ring&Simple&Left stabilizer&Right stabilizer\\\midrule
\ring{445}&3&$\{0,9,10,11\}\cong V_4$&$\{0,1,2,11\}\cong C_4$\\
\ring{493}&5&$\{0,9,10,11\}\cong V_4$&$\{0,3,4,11\}\cong C_4$\\\bottomrule\end{tabular}\end{center}
These groups are not isomorphic, so neither ring admits a complex categorification.
All subgroup products, element orders and grading products are checked directly.

The remaining four originally unresolved rings in this family are modeled:
$\ring{541}$ by a classical representation category, and
$\ring{428},\ring{513},\ring{597}$ by explicit twisted subgroup bimodules in
$\Vec_{A_4\times C_2}^{\omega}$ with a diagonal subgroup of order two.
Two invariant relative trilinear cocycles give a unified construction; their complete
values are retained, and every cocycle and subgroup-algebra identity is checked.
The independent binary-tetrahedral realization of $\ring{597}$ is described below.
Thus all six cases in the dimension-twenty-four family are settled.
